% ===== §7 Poeticity does not guarantee the good =====
\section{Poeticity Does Not Guarantee the Good}
\label{p09:sec:turn}

The preceding section raised the poem to a great height---the paradigm of the good sign, the Archimedean point of the semiotics. Such praise, if it is not at once and severely qualified, would bring the whole paper down here. For an objection waits at the door, and it is sharp enough to be fatal: \emb{poeticity guarantees only the self-continuation of generativity, not the good.}

The objection is fatal because it turns the paper's own blade against the paper. The \emph{antithesis} of this paper (\xref{p09:sec:antithesis}) laid down an iron law: \emb{self-continuation, in itself, is morally indifferent.} Capital's accumulation continues itself; the rose's blooming continues itself; the parasitic bond continues itself---``ceaseless circulation,'' alone, never constitutes the good. We must now concede: \emph{the poem is no exception.} A generator of $\lambda>1$, bearing a phase, unfurling without end---by what right is what it generates the good circle? By what right does the poetic spiral necessarily climb, rather than descend?

To speak honestly: there is no such guarantee. And the counter-examples lie ready to hand, and must be faced.

\subsection{The vicious poem exists}
\label{p09:sec:vicious-poem}

\begin{claim}[Poeticity does not entail the good]
There exist genuinely poetic, yet vicious, generators of generativity. Fascist aesthetics is poetic---it is by no means a dead grammar; it installs in its followers a grammar of meaning able to go on self-generating, self-``sublimating,'' impassioned, unending, demanding personal witness, refusing to be settled. Cultic speech is poetic. The addict's inner monologue is poetic---it ceaselessly generates new rationalizations, new cravings, self-reproducing, never closing. All satisfy the criterion of poeticity given in the preceding section: $\lambda>1$, the interpreter personally drawn in, refusal of settling, spiral and not mere repetition. They are \emph{genuine} generators of generativity, and yet vicious. Therefore ``poeticity $\Rightarrow$ good'' is a false proposition.
\end{claim}

We must pause and see exactly where the preceding section erred---for the error is the very one the paper had earlier warned itself to guard against. In \xref{p09:sec:synthesis} we set down a candid caveat: the correspondence ``sublimation $=$ positive phase'' requires an argument independent of definition, on pain of circularity. And \xref{p09:sec:semiotics}, in praising the poem, did precisely this---it quietly stitched ``poeticity'' to ``the good'' by the sign of the holonomy, without ever arguing, independently, \emph{why the poetic phase must be positive}. This section comes to discharge that debt.

The manner of discharge is to concede a stratification:

\begin{claim}[Poeticity is necessary but not sufficient for the good]
Poeticity guarantees the \emph{form} of the cycle---that it is a living spiral, not a dead circle (ossification) nor a settled closure (the contract). Poeticity guarantees that the cycle \emph{can} endure, \emph{can} sublimate. But it does \emph{not} determine the \emph{direction} of that sublimation. Poeticity gives the cycle the capacity for endurance, but not the direction of endurance. The good must be guaranteed by something else.
\end{claim}

\subsection{Two orthogonal criteria: poeticity and justice}
\label{p09:sec:orthogonal}

What is that ``something else''? To pursue it is to drive the paper's two languages to the point at which they must separate, and then be rejoined. I claim:

\begin{claim}[The orthogonality of poeticity and justice]
\emph{Poeticity} is a \emph{dynamical / formal} criterion: it asks---can this cycle endure? is it a spiral or has it gone rigid? \emph{Goodness} is a \emph{justice} criterion: it asks---is the phase of this cycle drawn from an internal reflux, or extracted from an outside? whose conditions of generation does it reproduce, and whom does it reduce to expendable fuel? The two criteria are \emph{orthogonal}: a cycle may be poetic and vicious (fascist aesthetics), or just and rigid (a bond fair but long since lifeless). The formal ``aliveness'' and the just ``goodness'' are two independent dimensions.
\end{claim}

With this orthogonal decomposition, fascist aesthetics can be diagnosed precisely, and no longer merely reviled as ``bad.'' It \emph{is} poetic---it meets every formal condition; it is \emph{not} just---its spiral can go on sublimating only because it takes ``the enemy,'' ``the alien,'' ``the other'' as expendable fuel. Its inner ``sublimation'' is \emb{stolen}: the phase is produced not by the tension among the participants internal to the cycle but extracted from an outside excluded and consumed, and then disguised as an internal sublimation.

This drives an earlier insight to its terminus. In \xref{p09:sec:antithesis} we said that the vicious cycle is ``an open system masquerading as a closed loop''---it dumps its entropy on an outside, deferring its own collapse by an ever-widening extraction. We now see that this structure has, at the level of poeticity, its exact twin:

\begin{claim}[The vicious poem is the stolen phase]
\emph{Fascist aesthetics is to the poem as capital is to generativity}---formal twins, ethical opposites. Capital ``self-augments,'' formally resembling generativity, but its increment is an extracted, counterfeit increment; the vicious poem ``self-sublimates,'' formally resembling the true poem, but its phase is a stolen, counterfeit phase. Both counterfeit a remainder that ought to have been internal, by excluding certain beings as pure fuel. A cycle's poeticity answers ``can it survive''; its justice answers ``\emph{should} it survive, and \emph{on whose blood} does it survive.''
\end{claim}

\subsection{An interrogable criterion: who is the pure fuel?}
\label{p09:sec:fuel}

The orthogonal decomposition gives us a concrete, reiterable diagnostic question---one that will become the point of departure for the next section's constructive work. How to tell the internal sublimation from the stolen one? Look whether the cycle has, within it, an outside systematically reduced to \emb{pure fuel}---a being that only contributes and never benefits: a silenced partner, a colonized other, a group consumed as an enemy, a future overdrawn, a nature exhausted.

\begin{claim}[The pure-fuel criterion]
Whether a spiral's sublimation is good or vicious may be interrogated by this question: \emph{is there, within this cycle, some one person or thing that exists only as cost---excluded from benefiting in the very sublimation it feeds?} If there is, this is a vampiric poeticity---its phase is stolen. In the good poetic cycle there is no ``pure outside'': every being drawn in is both producer and beneficiary of the phase; the sublimation refluxes to every one who generates it.
\end{claim}

This criterion is the ethical redemption of the ``open system masquerading as a closed loop'' of \xref{p09:sec:antithesis}, and it carries the core principle of generative justice---that value, unalienated, should return to those who create it---over from the domain of matter and labour into the domain of sublimation and witness. It will be locked, in the next section, into the substantive content of the ``appropriate'' poem.

It must finally be stressed where this section stands in the logic of the whole. To concede that ``poeticity does not guarantee the good'' does not weaken the preceding section; it is the very precondition of that section's standing. For it hands ``the good'' back to an independent, just criterion, and thereby forbids the poem---and the geometric form behind it---to usurp the role of a meta-framework that subsumes all the rest. The poetic framework can adjudicate ``sustainability / endurance''; but to adjudicate ``the good'' it must hand the microphone to the framework of justice. \emb{The poem answers whether the cycle can live; justice answers whether it should, and on whose blood.} These two questions are orthogonal, and therefore require two mutually unsubsuming frameworks---the polyphonic structure the paper's final section will avow, here for the first time appearing in a form impossible to evade: that even the most beautiful of things cannot, alone, vouch for the good.
